\subsection{Hardware runs and submitted evidence}
Two run folders are included. Run 1 starts at a camera estimated Tag 3 range of approximately $1.53\,\mathrm{m}$ and records a final range of approximately $0.24\,\mathrm{m}$. The planner terminal state is \texttt{STOPPED\_SHORT}, about $4\,\mathrm{cm}$ outside the nominal $0.20\,\mathrm{m}$ guard. Run 2 starts at approximately $1.58\,\mathrm{m}$ and reaches a recorded minimum of $0.296\,\mathrm{m}$ during \texttt{CHASING}. Its log later shows target loss and does not record \texttt{ARRIVED}. The student reports that the phone video shows the robot stop in front of Tag 3; that phone video should accompany the submission. ToF returned zero on all channels, so both controllers relied on camera range.

\begin{table}[ht]
\centering
\caption{Summary of the two saved hardware runs.}
\begin{tabular}{ll}
\toprule
Evidence & Recorded result \\
\midrule
Run 1 start / minimum range & $1.53 / 0.24\,\mathrm{m}$ \\
Run 1 terminal state & \texttt{STOPPED\_SHORT} \\
Run 2 start / minimum range & $1.58 / 0.296\,\mathrm{m}$ \\
Run 2 final logged state & \texttt{DISARMED\_NO\_ACK} after target loss \\
ToF & unavailable (zero readings) \\
Voltage basis & model estimates from wheel commands \\
\bottomrule
\end{tabular}
\end{table}

The second camera video and telemetry are in the Run 2 folder. The supplied screen recording shows a panel view during the session. The phone video of the mission is available at \url{https://youtu.be/K2ukuONWcgU}. The student reports that the phone video shows the robot stopping in front of Tag 3; the run logs retain their independently recorded range and state. The two CSV data sets remain separate.

For Run 2, position progress is computed from camera range as $d_0-d(t)$ rather than independent Cartesian odometry. Over the logged \texttt{CHASING} samples, range progress error has RMSE $0.0231\,\mathrm{m}$ and maximum absolute value $0.1368\,\mathrm{m}$. Velocity error has RMSE $0.0538\,\mathrm{m/s}$ and maximum absolute value $0.1526\,\mathrm{m/s}$, comparing the cubic reference with SDK chassis velocity.

\begin{figure}[ht]
\centering
\includegraphics[width=0.88\linewidth]{../data/hardware_20261118_011556_504834/attempt_01/analysis/references.png}
\caption{Run 2 reference and response. Position response uses camera range derived progress.}
\end{figure}
\begin{figure}[ht]
\centering
\includegraphics[width=0.88\linewidth]{../data/hardware_20261118_011556_504834/attempt_01/analysis/tracking_errors.png}
\caption{Run 2 reference minus response.}
\end{figure}

\begin{longtable}{@{}rrrrrl@{}}
\caption{One second motor voltage model samples from the recorded hardware run.}\label{tab:voltage}\\
\toprule
$t$ (s)&$V_1$ (V)&$V_2$ (V)&$V_3$ (V)&$V_4$ (V)&State\\\midrule
\endfirsthead
\toprule
$t$ (s)&$V_1$ (V)&$V_2$ (V)&$V_3$ (V)&$V_4$ (V)&State\\\midrule
\endhead
0&0.000&0.000&0.000&0.000&DISARMED\\
1&0.000&0.000&0.000&0.000&DISARMED\\
2&0.000&0.000&0.000&0.000&DISARMED\\
3&0.000&0.000&0.000&0.000&DISARMED\\
4&0.000&0.000&0.000&0.000&DISARMED\\
\bottomrule
\end{longtable}

The complete motor voltage model and high rate execution logs are supplied as \texttt{motor\_voltage\_log.csv} and \texttt{execution.csv}. Values are estimates from the assessment motor model, not direct voltage measurements.
